\section{Complete reconstruction-departure evaluation}
\label{app:reconstruction-results}

This appendix gives the numerical observation protocol, full method roster,
descriptive condition results, stability accounting, selection frequencies
and measured resources for Section~\ref{sec:reconstruction-mechanism}.
The frozen implementation and complete arrays are under
\path{extensions/reconstruction_mechanism}. All displays are generated
from the completed independent analysis; their source and output hashes are
recorded in
\path{generated/reconstruction_mechanism/ASSET_PROVENANCE.json}.

\subsection{Observation laws and separation of roles}

The linear negative control is
\[
 y_t=0.65y_{t-1}-0.15y_{t-2}+\epsilon_t,
 \qquad \epsilon_t\stackrel{\mathrm{iid}}\sim\mathcal N(0,0.1^2),
\]
with two initial zeros, 256 discarded stored values and 72 retained values
beginning at $y_{256}$. The nonlinear AR uses
\[
 y_t=\tanh\{1.25y_{t-1}-0.65y_{t-2}
             +0.25\sin(3y_{t-1})+u_t\},
 \qquad u_t\stackrel{\mathrm{iid}}\sim U[-0.1,0.1],
\]
two initial values uniform on $[-0.8,0.8]$, and 80 discarded stored values.
Van der Pol uses $\dot x=v$, $\dot v=(1-x^2)v-x$; Duffing uses
$\dot x=v$, $\dot v=-0.2v+x-x^3+0.3\cos(1.2t+\varphi)$.
The inherited float64 RK4 implementation has step 0.025, observation interval
0.1 and burn time ten. Initial $x,v$ are independent uniform on $[-2,2]$;
$\varphi$ is uniform on $[0,2\pi]$. Oscillator observations are
$0.95\tanh(x)+U[-0.02,0.02]$. The numerical integrator and observation
map define the tested laws; they are not exact continuous-ODE trajectories.

For each law and repetition, all split episodes are generated once. Additional
noise is $\eta s_A Z$, with $\eta\in\{0,0.1\}$,
$Z\stackrel{\mathrm{iid}}\sim\mathcal N(0,1)$ and $s_A$ the population
standard deviation of all unaugmented observed A episode values. Underlying
episodes and $Z$ arrays are paired across the two levels. The noise changes
both input histories and forecast targets. At zero added noise, the inherited
oscillator sensor noise remains. On B/C/D, the augmented linear observations are Gaussian conditional
on the realized A-estimated amplitude; across repetitions that random
amplitude induces a variance mixture. Neither exact stationarity after
finite burn-in nor a known-amplitude oracle interpretation of fitted ridge
is assumed. The Gaussian additions also preclude carrying over a bounded-loss
claim from the earlier bounded observation studies.

The 24 repetition seeds are 946100--946123. Development seed 975001 is
excluded. A/B/C/D base episodes and unit-noise arrays are generated from
distinct independent random streams; the added-noise amplitude is estimated
from A and shared across splits. Row caps and fitting seeds are paired between noise
levels. All 13 within-episode origins, at indices $11,15,\ldots,59$, use only
history through their origin and target offsets $(1,4,12)$. There are
\RMTargetCells{} scored target cells per method across the complete roster,
but only 24 independent aggregates in the primary inference. No individual
origin, horizon or paired-noise task is promoted to an independent replicate.

A fits every scaler, reference bank, kernel factor, raw predictor and decoder.
B selects model settings separately for each horizon. C selects each
non-ridge mixing coefficient, then one whole non-KAHM procedure. Selected
states are serialized before reserved queries are passed to prediction, and
targets enter the execution wrapper only after forecasts have been fixed.
No parameter or selection is updated on D.

\subsection{Fixed settings and complete descriptive losses}

The common locality logit is
$-\|z-\mu_j\|^2/(2s_j)$, where $s_j$ is the inherited mean squared
local radius, floored at $10^{-6}$. Softmax normalization occurs separately
within the 2-, 4- and 8-model blocks. Principal rank is two, there are at most
24 references per local model, and the A bank cap is 512 rows. The inherited
KAHM reference-energy regularization is unchanged. Reference and query kernels
use the same projected covariance precision; the query log-kernel maximum is
subtracted before exponentiation for both KAHM and positive RBF ratios.

The inherited history scaler replaces standard deviations at or below
$10^{-12}$ by one. The subsequent A feature scalers use
$\max(\mathrm{sd},10^{-6})$. Raw ridge/RBF history scalers also use this
$10^{-6}$ floor. Common and departure blocks are radially bounded separately,
with no scaling after the transform. The unbounded ablations retain the same
bounded common block and omit only the departure radial map. Their failures
or successes cannot replace the bounded primary candidate.

Every ridge head uses the mean-SSE objective and the five penalties
$10^{-4},10^{-3},10^{-2},10^{-1},1$. The equivalent sum-SSE penalty is
$n\alpha$, not $\alpha$. The raw RBF arm concatenates standardized histories
with at most 64 Nystr\"om coordinates, then scales those concatenated columns
using A alone. Bandwidth factors 0.5, 1 and 2 multiply the median positive
landmark squared-distance scale; B selects bandwidth and penalty separately
by horizon. Histogram gradient boosting uses 100 iterations, L2 penalty one,
learning rates 0.05 or 0.1 and maximum leaves 7 or 15; B again selects separately
by horizon. Persistence repeats the latest observed value before its common
baseline mixing step. All non-ridge procedures use the same C mixing grid
$\{0,1/4,1/2,1\}$, with deterministic tie order.

\begin{table}[!htbp]
\centering\scriptsize
\setlength{\tabcolsep}{3pt}
\caption{All ten procedures, every declared law and both added-noise levels.
Entries are mean horizon-averaged A-variance-normalized D MSE over 24
repetitions; lower is better. The first six columns enter the nonlinear
aggregate with equal weight. The final two form the separate linear control.
These conditionwise summaries are descriptive. ``Unbounded'' omits only the
departure radial bound. Persistence and raw nonlinear controls include their
C-selected blend with ridge.}
\label{tab:mechanism-all-methods}
% Generated from the complete saved reconstruction-mechanism analysis.
\begin{tabular}{@{}lrrrrrrrr@{}}
\toprule
Procedure & \multicolumn{2}{c}{Nonlinear AR} & \multicolumn{2}{c}{Van der Pol} & \multicolumn{2}{c}{Duffing} & \multicolumn{2}{c}{Linear AR} \\
\midrule
Added noise $\eta$ & 0 & 0.1 & 0 & 0.1 & 0 & 0.1 & 0 & 0.1 \\
Raw ridge & 0.0216 & 0.0344 & 0.0750 & 0.1361 & 0.0192 & 0.0681 & 0.8951 & 0.8975 \\
Raw RBF & 0.0197 & 0.0323 & 0.0131 & 0.0647 & 0.0129 & 0.0442 & 0.8969 & 0.8994 \\
Boosted trees & 0.0201 & 0.0323 & 0.0040 & 0.0586 & 0.0133 & 0.0481 & 0.8975 & 0.8997 \\
Persistence blend & 0.0219 & 0.0345 & 0.0750 & 0.1362 & 0.0192 & 0.0684 & 0.8955 & 0.8988 \\
Common basis & 0.0190 & 0.0310 & 0.0251 & 0.0686 & 0.0157 & 0.0460 & 0.8959 & 0.8983 \\
KAHM departure (bounded) & 0.0192 & 0.0314 & 0.0101 & 0.0622 & 0.0129 & 0.0439 & 0.8972 & 0.8990 \\
RBF departure (bounded) & 0.0189 & 0.0309 & 0.0115 & 0.0634 & 0.0133 & 0.0436 & 0.8968 & 0.8989 \\
KAHM departure (unbounded) & 0.0198 & 0.0321 & 0.0108 & 0.0619 & 0.0129 & 0.0440 & 0.8975 & 0.8995 \\
RBF departure (unbounded) & 0.0190 & 0.0310 & 0.0118 & 0.0627 & 0.0127 & 0.0437 & 0.8972 & 0.8993 \\
C-selected non-KAHM policy & 0.0192 & 0.0312 & 0.0040 & 0.0590 & 0.0127 & 0.0434 & 0.8978 & 0.8997 \\
\bottomrule
\end{tabular}

\end{table}

The machine-readable files \texttt{horizon\_metrics.csv} and
\texttt{condition\_horizon\_metrics.csv} retain every method, law, noise
level and horizon, including raw MSE and the A-only loss scale. They are not
pooled with historical losses that use different normalizations. The
\texttt{replicate\_contrasts.csv} file contains all 24 margin-gap values for
each primary contrast, allowing the six lower bounds to be recomputed without
importing any forecasting model.

\subsection{Stability, numerical safeguards and correction activation}

The bounded-feature checks cover 192 tasks for each of three bounded
procedures. Coefficient and correction checks additionally distinguish the
three horizons, giving 576 checks per procedure for each inequality. The
independent analyzer checks the recorded fitted-state bound formulas and
recomputes actual D corrections as the saved forecast minus the saved raw
ridge forecast. It does not claim to reconstruct the fitted coefficient
matrices from summary statistics alone. Table~\ref{tab:mechanism-bounds}
reports maxima over the complete roster; maxima in different columns need
not come from the same task or horizon.

\begin{table}[!htbp]
\centering\small
\setlength{\tabcolsep}{5pt}
\caption{Observed ratios to the deterministic bounded-feature and
ridge-readout limits. Feature ratios concern the common block for the
common procedure and the respective departure block for the two appended
procedures; the readout bound uses the complete concatenated feature radius.
A zero correction with a zero bound contributes ratio zero. ``Violations''
sums failed feature, coefficient and correction checks at the fixed numerical
tolerances. These are implementation checks, not confidence intervals for
forecast risk.}
\label{tab:mechanism-bounds}
% Generated from the complete saved reconstruction-mechanism analysis.
\begin{tabular}{@{}lrrrr@{}}
\toprule
Bounded procedure & \shortstack{Feature norm\\/ radius} & \shortstack{Coefficient norm\\/ bound} & \shortstack{Correction\\/ bound} & Violations \\
\midrule
Common basis & 0.99493 & 0.56103 & 0.29634 & 0 \\
KAHM departure (bounded) & 0.99943 & 0.55231 & 0.20110 & 0 \\
RBF departure (bounded) & 0.99734 & 0.44963 & 0.19674 & 0 \\
\bottomrule
\end{tabular}

\end{table}

The maximum taskwise ratio of KAHM maximum absolute error to ridge maximum
absolute error is \RMMaxErrorRatio{}, and \RMStabilityFailures{} tasks fail
at least one declared stability screen. Local KAHM reconstruction fallback
affects \RMKFallbackABC{} A/B/C rows and \RMKFallbackD{} D rows across all
tasks; the corresponding RBF D count is \RMRFallbackD{}. Counts are unions
over local models within a task and split, so the same row is not counted
again merely because several local maps invoke their guard. The fixed local
fallback sets the affected departure to zero. The complete valid analysis
requires zero wholly failed methods and zero final prediction-cell recoveries;
neither is hidden by removing a task or substituting a successful subset.

\begin{table}[!htbp]
\centering\small
\setlength{\tabcolsep}{4pt}
\caption{C-selected correction activation. Nonlinear and linear counts have
432 and 144 task/horizon cells, respectively. The four mixing-weight columns
combine all 576 cells. A policy choice of raw ridge is counted as zero
correction; raw ridge itself has no mixing selection. Nonzero mixing is not
an empirical guarantee of D improvement, and no minimum activation frequency
was part of the decision.}
\label{tab:mechanism-activation}
% Generated from the complete saved reconstruction-mechanism analysis.
\begin{tabular}{@{}lrrrrrr@{}}
\toprule
Procedure & \shortstack{Nonlinear\\active / 432} & \shortstack{Linear\\active / 144} & $\gamma=0$ & $\gamma=1/4$ & $\gamma=1/2$ & $\gamma=1$ \\
\midrule
Raw ridge & --- & --- & --- & --- & --- & --- \\
Raw RBF & 427 & 56 & 93 & 52 & 158 & 273 \\
Boosted trees & 396 & 48 & 132 & 106 & 175 & 163 \\
Persistence blend & 70 & 32 & 474 & 97 & 5 & 0 \\
Common basis & 430 & 56 & 90 & 9 & 66 & 411 \\
KAHM departure (bounded) & 432 & 54 & 90 & 23 & 79 & 384 \\
RBF departure (bounded) & 431 & 58 & 87 & 14 & 71 & 404 \\
KAHM departure (unbounded) & 428 & 51 & 97 & 40 & 107 & 332 \\
RBF departure (unbounded) & 429 & 61 & 86 & 26 & 75 & 389 \\
C-selected non-KAHM policy & 430 & 63 & 83 & 46 & 100 & 347 \\
\bottomrule
\end{tabular}

\end{table}

\begin{table}[!htbp]
\centering\small
\setlength{\tabcolsep}{8pt}
\caption{Whole-procedure choices made by the C-selected non-KAHM policy.
Each task contributes one choice across all horizons after each candidate's
own horizon-wise C mixing selection. Both KAHM arms are excluded. These are
validation choices, not choices made from D errors.}
\label{tab:mechanism-policy}
% Generated from the complete saved reconstruction-mechanism analysis.
\begin{tabular}{@{}lrr@{}}
\toprule
C-selected complete procedure & Nonlinear / 144 & Linear / 48 \\
\midrule
Raw ridge & 0 & 1 \\
Raw RBF & 19 & 6 \\
Boosted trees & 51 & 8 \\
Persistence blend & 2 & 20 \\
Common basis & 15 & 4 \\
RBF departure (bounded) & 19 & 1 \\
RBF departure (unbounded) & 38 & 8 \\
\bottomrule
\end{tabular}

\end{table}

\subsection{Measured costs and reproducibility limits}

\begin{table}[!htbp]
\centering\small
\setlength{\tabcolsep}{4pt}
\caption{Mean charged resources over all 192 equal-sized task suites.
Fit/select includes required A fitting, B selection and C calibration.
D prediction covers all 1,664 histories and three horizons. Shared computation
nodes are charged once to every procedure requiring them; policy fitting pays
the union of its non-KAHM candidates, while its prediction/state follow the
selected complete procedure. State denotes serialized selected objects;
arrays count their owned NumPy buffers. These are measured implementation
costs, not minimum deployment complexity.}
\label{tab:mechanism-resources}
% Generated from the complete saved reconstruction-mechanism analysis.
\begin{tabular}{@{}lrrrrr@{}}
\toprule
Procedure & \shortstack{Fit/select\\wall (s)} & \shortstack{Fit/select\\CPU (s)} & \shortstack{D predict\\wall (ms)} & \shortstack{State\\(KiB)} & \shortstack{Arrays\\(KiB)} \\
\midrule
Raw ridge & 0.0010 & 0.0010 & 0.147 & 1.10 & 0.63 \\
Raw RBF & 0.0135 & 0.0135 & 2.925 & 65.58 & 64.28 \\
Boosted trees & 0.6452 & 0.6452 & 19.046 & 373.94 & 335.61 \\
Persistence blend & 0.0011 & 0.0011 & 0.197 & 1.15 & 0.66 \\
Common basis & 0.0367 & 0.0353 & 5.363 & 9.00 & 6.82 \\
KAHM departure (bounded) & 0.0833 & 0.0819 & 29.245 & 118.81 & 114.76 \\
RBF departure (bounded) & 0.0658 & 0.0644 & 25.758 & 55.32 & 51.85 \\
KAHM departure (unbounded) & 0.0812 & 0.0798 & 26.655 & 118.81 & 114.76 \\
RBF departure (unbounded) & 0.0642 & 0.0628 & 23.810 & 55.32 & 51.85 \\
C-selected non-KAHM policy & 0.7374 & 0.7360 & 15.068 & 151.36 & 137.94 \\
\bottomrule
\end{tabular}

\end{table}

Each task runs serially in a fresh process with one assigned CPU and one
numerical thread. The charged primitive-node costs exclude diagnostic-only
rank calculations; full-suite timings, including diagnostics, are retained
separately. Node accounting includes model candidates that are actually
evaluated, rather than only the winning hyperparameters. Wall and CPU times
depend on the recorded machine and software environment. The peak whole-task
process RSS is \RMPeakProcessMiB{} MiB; it includes imports, all ten arms,
shared diagnostics and evidence arrays, and is not a per-method RAM estimate.
Serialized state and owned array buffers omit transient workspace and do not
measure online update latency. No cost ceiling was imposed without an
application-derived requirement.

The independent analyzer verifies the immutable local source/input freeze,
the complete task roster, exact target correspondence, A variance scales,
policy predictions, nonfinite/failure accounting and all primary aggregation.
The reporting builder separately reconstructs condition and repetition means
and all six $t$ lower bounds before producing these displays. Numerical
development checks, model replay and independent calculations are consistency
checks; they do not constitute a new scientific replication. The protocol is
locally prospective, informed by previous results, and explicitly not public
preregistration. Neither these known-law repetitions nor the algebraic
distinctness certificate establishes a natural-application forecasting claim.
